% ═══════════════════════════════════════════════════════════════
% LH-01b: Mechanische Energie – LEHRERHINWEISE
% Physik Klasse 10 MR
% ═══════════════════════════════════════════════════════════════

\documentclass[11pt, a4paper]{article}

\usepackage[utf8]{inputenc}
\usepackage[T1]{fontenc}
\usepackage[ngerman]{babel}

\usepackage{helvet}
\renewcommand{\familydefault}{\sfdefault}

\usepackage{microtype}
\usepackage[margin=2cm]{geometry}
\usepackage{enumitem}
\usepackage{tabularx}
\usepackage{booktabs}
\usepackage{array}
\usepackage{longtable}

\usepackage{amsmath, amssymb}
\usepackage{siunitx}
\sisetup{locale = DE, per-mode = symbol}

\usepackage{fancyhdr}
\usepackage{lastpage}

\usepackage{xcolor}
\usepackage{tcolorbox}
\tcbuselibrary{skins}

% ═══════════════════════════════════════════════════════════════
% FARBEN
% ═══════════════════════════════════════════════════════════════

\definecolor{physik}{HTML}{2c5aa0}
\definecolor{hellgrau}{HTML}{F5F5F5}
\definecolor{mittelgrau}{HTML}{888888}
\definecolor{loesung}{HTML}{c62828}
\definecolor{hinweis}{HTML}{b35c00}

\colorlet{fachfarbe}{physik}

% ═══════════════════════════════════════════════════════════════
% VARIABLEN
% ═══════════════════════════════════════════════════════════════

\newcommand{\fach}{Physik}
\newcommand{\thema}{Mechanische Energie}
\newcommand{\klasse}{10 MR}
\newcommand{\stession}{01b}
\newcommand{\dauer}{90 Min (DS)}

% ═══════════════════════════════════════════════════════════════
% KOPF- UND FUSSZEILE
% ═══════════════════════════════════════════════════════════════

\pagestyle{fancy}
\fancyhf{}
\renewcommand{\headrulewidth}{0.4pt}
\renewcommand{\footrulewidth}{0pt}
\lhead{\small \fach{} Klasse \klasse{} | \thema{} | \textbf{Lehrerhinweise}}
\rhead{\small \dauer}
\cfoot{\small\textcolor{mittelgrau}{Seite \thepage{} von \pageref{LastPage}}}

% ═══════════════════════════════════════════════════════════════
% BOXEN
% ═══════════════════════════════════════════════════════════════

\newtcolorbox{infobox}[1][]{
    colback=hellgrau,
    colframe=fachfarbe,
    coltitle=white,
    colbacktitle=fachfarbe,
    boxrule=1pt,
    arc=0pt,
    left=8pt, right=8pt, top=6pt, bottom=6pt,
    before skip=8pt, after skip=8pt,
    fonttitle=\bfseries,
    title=#1
}

\newtcolorbox{warnbox}[1][]{
    colback=white,
    colframe=hinweis,
    coltitle=white,
    colbacktitle=hinweis,
    boxrule=1pt,
    arc=0pt,
    left=8pt, right=8pt, top=6pt, bottom=6pt,
    before skip=8pt, after skip=8pt,
    fonttitle=\bfseries,
    title=#1
}

% ═══════════════════════════════════════════════════════════════
% BEFEHLE
% ═══════════════════════════════════════════════════════════════

\newcommand{\lsg}[1]{\textcolor{loesung}{\textbf{#1}}}

\setlength{\parindent}{0pt}
\setlength{\parskip}{6pt}

% ═══════════════════════════════════════════════════════════════
% DOKUMENT
% ═══════════════════════════════════════════════════════════════

\begin{document}

% ═══════════════════════════════════════════════════════════════
% HEADER
% ═══════════════════════════════════════════════════════════════

\begin{center}
    {\Large\bfseries\textcolor{fachfarbe}{Lehrerhinweise: \thema}}\\[0.3em]
    {\small \fach{} Klasse \klasse{} | Stunde \stession{} | \dauer}
\end{center}
\vspace{-0.3em}
\hrule
\vspace{0.8em}

% ═══════════════════════════════════════════════════════════════
% ÜBERSICHT
% ═══════════════════════════════════════════════════════════════

\section*{Übersicht}

\begin{tabularx}{\textwidth}{@{}lX@{}}
\textbf{Rahmenplan:} & RP RS/GES 2021, Kl. 10: Energie – Energiegewinnung \\
\textbf{Vorwissen:} & $E_\text{pot}$, $E_\text{kin}$ aus Klasse 7 (Wiederholung) \\
\textbf{Materialien:} & LP-01b (digital), AB-01b (Print), MT-01b (digital), Beamer \\
\textbf{BE gesamt:} & AB: 24 BE, MT: 20 BE \\
\end{tabularx}

\begin{infobox}[Lernziele (Feinziele)]
\begin{itemize}[leftmargin=2em, itemsep=2pt]
    \item FZ1 (AFB I): SuS nennen die Definition von Energie
    \item FZ2 (AFB I): SuS geben die Formel $E_\text{pot} = m \cdot g \cdot h$ an
    \item FZ3 (AFB I): SuS geben die Formel $E_\text{kin} = \frac{1}{2} m v^2$ an
    \item FZ4 (AFB II): SuS berechnen $E_\text{pot}$ und $E_\text{kin}$
    \item FZ5 (AFB II): SuS wenden den Energieerhaltungssatz an
    \item FZ6 (AFB II): SuS zeichnen Energieflussdiagramme (ideal + real)
    \item FZ7 (AFB III): SuS erklären Energieumwandlungen (ideal + real mit Reibung)
    \item FZ8 (AFB III): SuS berechnen $v$ aus $h$ bzw. $h$ aus $v$ (Transfer)
\end{itemize}
\end{infobox}

% ═══════════════════════════════════════════════════════════════
% STUNDENVERLAUF
% ═══════════════════════════════════════════════════════════════

\section*{Stundenverlauf}

\subsection*{Teil 1: Lernpfad + Arbeitsblatt (0–70')}

\begin{tabularx}{\textwidth}{@{}c|l|X|c@{}}
\toprule
\textbf{Zeit} & \textbf{LP-Schritt} & \textbf{Inhalt} & \textbf{AB} \\
\midrule
0–5' & Schritt 1 & Was hat das mit mir zu tun? Motivationseinstieg & – \\
5–10' & Schritt 2 & Was ist Energie? Definition, Einheit J & K1 \\
10–17' & Schritt 3 & $E_\text{pot} = m \cdot g \cdot h$ einführen & K2 \\
17–25' & Schritt 4 & Simulation: Lageenergie, Übung 1 & A1a \\
25–32' & Schritt 5 & $E_\text{kin} = \frac{1}{2} m v^2$ einführen & K3 \\
32–40' & Schritt 6 & Simulation: Bewegungsenergie, Übung 2 & A1b \\
40–45' & Schritt 7 & Energieerhaltungssatz & K4 \\
45–52' & Schritt 8 & Simulation: Rampe + Energieumwandlung & A2 \\
52–58' & Schritt 9 & Reales Verhalten (Reibung → Wärme) & A3 \\
58–68' & Schritt 10 & Knobelaufgaben (Transfer: v aus h, h aus v) & A4 \\
68–70' & Schritt 11 & Zusammenfassung, AB-Check & – \\
\bottomrule
\end{tabularx}

\vspace{0.5em}
\textit{K = Kasten, A = Aufgabe auf dem AB}

\subsection*{Teil 2: Stundenleistung (70–90')}

\begin{tabularx}{\textwidth}{@{}c|l|X@{}}
\toprule
\textbf{Zeit} & \textbf{Phase} & \textbf{Inhalt} \\
\midrule
70–72' & Vorbereitung & AB weglegen, MT-01b öffnen, Namen eingeben \\
72–90' & Bearbeitung & 6 Aufgaben, 20 BE, selbstständig \\
90' & Abschluss & Automatische Auswertung, Ergebnis notieren \\
\bottomrule
\end{tabularx}

\begin{warnbox}[Zeitpuffer]
Bei Zeitmangel: AB Aufgabe 3 (Energieflussdiagramm) kürzen oder als HA.\\
Bei Zeitüberschuss: Zusatzfrage „Was passiert mit der Energie beim Bremsen?"
\end{warnbox}

% ═══════════════════════════════════════════════════════════════
% DIFFERENZIERUNG
% ═══════════════════════════════════════════════════════════════

\section*{Differenzierung}

\begin{tabularx}{\textwidth}{@{}c|l|X@{}}
\toprule
\textbf{Niveau} & \textbf{Aufgaben} & \textbf{Hinweise} \\
\midrule
$\star$ Basis & A1a & Nur $E_\text{pot}$ einsetzen, keine Umstellung \\
$\star\star$ Standard & A1b, A2, A3 & $E_\text{kin}$ berechnen (v quadrieren!), Erhaltung anwenden, Energieflussdiagramm \\
$\star\star\star$ Erweitert & A4 & Transfer-Aufgaben: v aus h bzw. h aus v durch Probieren ermitteln \\
\bottomrule
\end{tabularx}

% ═══════════════════════════════════════════════════════════════
% HÄUFIGE FEHLER
% ═══════════════════════════════════════════════════════════════

\section*{Häufige Schülerfehler}

\begin{tabularx}{\textwidth}{@{}l|X|X@{}}
\toprule
\textbf{Fehler} & \textbf{Diagnose} & \textbf{Reaktion} \\
\midrule
$g = 9{,}81$ verwenden & Falsch gelernter Wert & „Wir rechnen IMMER mit $g = 10$ m/s²!" \\
\midrule
v nicht quadrieren & Formel vergessen: $E_\text{kin} = \frac{1}{2} m v^2$ & „Bei $E_\text{kin}$ wird v hoch 2 genommen!" \\
\midrule
Einheit vergessen & Nur Zahl, kein J & „Energie hat IMMER die Einheit Joule (J)" \\
\midrule
$E_\text{kin} \neq E_\text{pot}$ beim Fallen & Energieerhaltung nicht verstanden & „Die gesamte $E_\text{pot}$ wird zu $E_\text{kin}$!" \\
\midrule
v in km/h verwenden & Keine Umrechnung & „v muss in m/s sein! 36 km/h = 10 m/s" \\
\bottomrule
\end{tabularx}

% ═══════════════════════════════════════════════════════════════
% LÖSUNGEN AB
% ═══════════════════════════════════════════════════════════════

\newpage
\section*{Lösungen zum Arbeitsblatt}

\subsection*{Merkkästen}

\textbf{Kasten 1:}
\begin{itemize}[leftmargin=2em]
    \item Energie = Fähigkeit, \lsg{Arbeit} zu verrichten oder \lsg{Wärme} abzugeben
    \item Einheit: \lsg{Joule (J)}
\end{itemize}

\textbf{Kasten 2:}
\begin{itemize}[leftmargin=2em]
    \item $E_\text{pot}$ in \lsg{J}
    \item $m$ in \lsg{kg}
    \item $g$ = \lsg{10} m/s²
    \item $h$ in \lsg{m}
\end{itemize}

\textbf{Kasten 3:}
\begin{itemize}[leftmargin=2em]
    \item $E_\text{kin}$ in \lsg{J}
    \item $m$ in \lsg{kg}
    \item $v$ in \lsg{m/s}
    \item v wird \lsg{quadriert}, doppelte v $\rightarrow$ \lsg{vierfache} Energie
\end{itemize}

\textbf{Kasten 4:}
\begin{itemize}[leftmargin=2em]
    \item Energie kann nur \lsg{umgewandelt} werden
\end{itemize}

\subsection*{Aufgaben}

\begin{tabularx}{\textwidth}{@{}l|X|c@{}}
\toprule
\textbf{Aufgabe} & \textbf{Lösung} & \textbf{BE} \\
\midrule
1a) $E_\text{pot}$ & $E_\text{pot} = 10\,\text{kg} \cdot 10\,\text{m/s}^2 \cdot 8\,\text{m} = \lsg{800\,\text{J}}$ & 2 \\
\midrule
1b) $E_\text{kin}$ & $E_\text{kin} = \frac{1}{2} \cdot 0{,}5\,\text{kg} \cdot (20\,\text{m/s})^2 = \lsg{100\,\text{J}}$ & 3 \\
\midrule
2a) $E_\text{pot}$ oben & $E_\text{pot} = 5\,\text{kg} \cdot 10\,\text{m/s}^2 \cdot 20\,\text{m} = \lsg{1000\,\text{J}}$ & 2 \\
\midrule
2b) $E_\text{kin}$ unten & $E_\text{kin} = \lsg{1000\,\text{J}}$ & 1 \\
\midrule
2c) Begründung & \lsg{Energieerhaltung: $E_\text{pot}$ wird vollständig zu $E_\text{kin}$} & 2 \\
\midrule
3a) Diagramm & $E_\text{pot} \rightarrow E_\text{kin} + E_\text{Wärme}$ (realistisch) & 2 \\
\midrule
3b) Erklärung & \lsg{$E_\text{pot}$ wird umgewandelt, durch Reibung geht Teil als Wärme verloren} & 4 \\
\midrule
4a) Auto h aus v & $E_\text{kin} = 50.000\,\text{J} \rightarrow h = \lsg{5\,\text{m}}$ & 4 \\
\midrule
4b) Stein v aus h & $E_\text{pot} = 100\,\text{J} \rightarrow v = \lsg{10\,\text{m/s}}$ & 4 \\
\bottomrule
\end{tabularx}

% ═══════════════════════════════════════════════════════════════
% LÖSUNGEN MT
% ═══════════════════════════════════════════════════════════════

\section*{Lösungen zur Stundenleistung (MT-01b)}

\begin{tabularx}{\textwidth}{@{}l|X|c@{}}
\toprule
\textbf{Aufgabe} & \textbf{Lösung} & \textbf{BE} \\
\midrule
1) MC & B: „Fähigkeit, Arbeit zu verrichten oder Wärme abzugeben" & 2 \\
\midrule
2) Zuordnung & $E \rightarrow$ J, $m \rightarrow$ kg, $v \rightarrow$ m/s, $g \rightarrow$ m/s² & 4 \\
\midrule
3) $E_\text{pot}$ & $E_\text{pot} = 4 \cdot 10 \cdot 5 = \lsg{200\,\text{J}}$ & 3 \\
\midrule
4) $E_\text{kin}$ & $E_\text{kin} = \frac{1}{2} \cdot 80 \cdot 25 = \lsg{1000\,\text{J}}$ & 3 \\
\midrule
5a) $E_\text{pot}$ oben & $E_\text{pot} = 2 \cdot 10 \cdot 10 = \lsg{200\,\text{J}}$ & 2 \\
\midrule
5b) $E_\text{kin}$ unten & $E_\text{kin} = \lsg{200\,\text{J}}$ (Energieerhaltung) & 2 \\
\midrule
6) Erklärung & Verwendung aller Begriffe, logische Erklärung & 4 \\
\bottomrule
\end{tabularx}

\vspace{0.5em}
\textbf{Notenschlüssel (6-Noten, MR Klasse 10):}

\begin{tabularx}{\textwidth}{@{}c|c|c|c|c|c|c@{}}
\toprule
\textbf{Note} & 1 & 2 & 3 & 4 & 5 & 6 \\
\midrule
\textbf{Prozent} & $\geq$96\% & $\geq$80\% & $\geq$60\% & $\geq$40\% & $\geq$20\% & $<$20\% \\
\midrule
\textbf{BE} & $\geq$19 & $\geq$16 & $\geq$12 & $\geq$8 & $\geq$4 & $<$4 \\
\bottomrule
\end{tabularx}

% ═══════════════════════════════════════════════════════════════
% MATERIAL-CHECKLISTE
% ═══════════════════════════════════════════════════════════════

\section*{Material-Checkliste}

\begin{itemize}[leftmargin=2em]
    \item[$\square$] Beamer + Laptop (für LP)
    \item[$\square$] AB-01b-energie (Klassensatz, 2-seitig)
    \item[$\square$] WLAN für SuS-Geräte (MT digital)
    \item[$\square$] LP-URL / QR-Code bereit
\end{itemize}

\begin{infobox}[URLs]
\textbf{Lernpfad:} \texttt{LP-01b-energie.html}\\
\textbf{Stundenleistung:} \texttt{MT-01b-energie.html}
\end{infobox}

% ═══════════════════════════════════════════════════════════════
% AUSBLICK
% ═══════════════════════════════════════════════════════════════

\section*{Ausblick}

\textbf{Nächste Stunde (01c):} Energiegewinnung, Kraftwerke, Wirkungsgrad

\textbf{Verknüpfung:} Mechanische Energie $\rightarrow$ Generatoren $\rightarrow$ elektrische Energie

\vfill
\hrule
\vspace{0.3em}
\begin{center}
\small\textcolor{mittelgrau}{LH-01b | Stand: Januar 2026}
\end{center}

\end{document}
